\section{증명해 볼 난제 후보}

2026-10-12. 아직 고르지 않음. 다룰 때는 (1) 증명한 부분, (2) 특수한
경우·부분 결과, (3) 막힌 지점으로 나눠 기록한다.

\begin{longtable}[]{@{}
  >{\raggedright\arraybackslash}p{(\columnwidth - 4\tabcolsep) * \real{0.3333}}
  >{\raggedright\arraybackslash}p{(\columnwidth - 4\tabcolsep) * \real{0.3333}}
  >{\raggedright\arraybackslash}p{(\columnwidth - 4\tabcolsep) * \real{0.3333}}@{}}
\toprule\noalign{}
\begin{minipage}[b]{\linewidth}\raggedright
문제
\end{minipage} & \begin{minipage}[b]{\linewidth}\raggedright
내용
\end{minipage} & \begin{minipage}[b]{\linewidth}\raggedright
현재 상태
\end{minipage} \\
\midrule\noalign{}
\endhead
\bottomrule\noalign{}
\endlastfoot
콜라츠 추측 & 짝수면 \(n/2\), 홀수면 \(3n+1\)을 반복하면 항상 1에
도달하는가 & 미해결 \\
에르되시--슈트라우스 추측 & \(n\ge2\)에서
\(\dfrac4n=\dfrac1x+\dfrac1y+\dfrac1z\)인 양의 정수가 항상 있는가 &
미해결. 작은 경우 계산·구성적 접근이 가능 \\
골드바흐 추측 & 2보다 큰 짝수는 두 소수의 합인가 & 미해결 \\
쌍둥이 소수 추측 & 차가 2인 소수 쌍이 무한히 많은가 & 미해결.
Zhang--Maynard 이후 간격 246 이하는 무한히 많음 \\
외로운 주자 추측 & 서로 다른 일정 속도로 원형 트랙을 도는 주자들 &
미해결 \\
하드비거--넬슨 문제 & 평면의 채색수 & 5, 6, 7 중 하나 \\
\end{longtable}

추천 출발점: 에르되시--슈트라우스, 외로운 주자 (작은 경우를 계산하거나
구성하기 쉬움).
